
\section{Nicholl-Lee-Nicholl Algorithm}

\subsection{History}
The Nicholl-Lee-Nicholl algorithm was published in 1987 by researchers Tina M. Nicholl, D.T. Lee, and Robin A. Nicholl. The NLN algorithm was developed specifically to address the severe computational inefficiencies found in earlier rectangular clipping algorithms \cite{MatthesDrakopoulos2022}.
In older methods like Cohen-Sutherland, algorithms would frequently calculate exact intersection points that were ultimately found to lie outside the clipping window, thus wasting valuable processing power \cite{NichollLeeNicholl1987}. The NLN method was designed to drastically reduce these redundant mathematical divisions and intersection calculations, rendering it one of the most highly optimized algorithms specifically for 2D rectangular clipping windows \cite{MatthesDrakopoulos2022}.

\subsection{General Description}
Unlike Cyrus-Beck, which can handle complex convex polygons, the NLN algorithm sacrifices shape versatility for extreme speed, strictly tailoring its logic for 2D upright rectangular clipping windows \cite{MatthesDrakopoulos2022}. It functions as a region-based algorithm but uses geometric space-division rather than simple endpoint encoding \cite{NichollLeeNicholl1987}.
The NLN algorithm processes line clipping through a highly structured sequence of geometric evaluations. It begins by locating the starting point ($P_0$) of the line segment relative to the clipping window. The 2D plane is conceptually divided into nine distinct regions: one inside the window, four edge or face regions, and four corner regions. $P_0$ is then assigned to one of these zones \cite{NichollLeeNicholl1987}. To reduce the number of conditional checks required by the computer processor, the NLN algorithm heavily utilizes geometric symmetry. Specific coordinate transformations are applied so that the algorithm only needs to process a subset of base cases, such as assuming $P_0$ is either strictly inside, located in a specific edge region, or located in a specific corner region \cite{NichollLeeNicholl1987}.
Once the starting region is identified, instead of infinitely extending the window's boundaries, the algorithm draws conceptual rays from $P_0$ through the four corners of the clipping window. These rays act as dividing lines that partition the remaining 2D space into highly specific sub-regions, which correspond to Left, Top, Right, or Bottom intersection zones. The algorithm then determines which of these sub-regions contains the second endpoint ($P_1$). Crucially, this determination is accomplished using slope comparisons or cross products rather than floating-point division, which saves a significant amount of computational time \cite{NichollLeeNicholl1987, MatthesDrakopoulos2022}. Because the plane has been systematically divided based on the precise geometry of $P_0$ and the window's corners, pinpointing the specific sub-region of $P_1$ immediately dictates exactly which window boundaries the line intersects. Consequently, an intersection equation is only executed if an intersection is absolutely guaranteed to lie on the visible boundary, entirely eliminating redundant mathematical operations \cite{MatthesDrakopoulos2022}.